% Part: first-order-logic
% Chapter: tableaux
% Section: identity

\documentclass[../../../include/open-logic-section]{subfiles}

\begin{document}

\olfileid{fol}{tab}{ide}

\olsection{\usetoken{P}{tableau} وچ \usetoken{S}{identity} دی شمولیت}

!!^{tableau}s وچ !!{identity} لئی وادھو استنباطی قاعدے چاہیدے نیں۔
$\eq$ دے قاعدے ایہ نیں ($t$، $t_1$ تے $t_2$ بند ٹرم نیں):

\begin{defish}
\AxiomC{}
\RightLabel{$\eq$}
\UnaryInfC{\sFmla{\True}{\eq[t][t]}}
\DisplayProof
\hfill
\AxiomC{\sFmla{\True}{\eq[t_1][t_2]}}
\noLine
\UnaryInfC{\sFmla{\True}{!A(t_1)}}
\RightLabel{$\TRule{\True}{\eq}$}
\UnaryInfC{\sFmla{\True}{!A(t_2)}}
\DisplayProof
\hfill
\AxiomC{\sFmla{\True}{\eq[t_1][t_2]}}
\noLine
\UnaryInfC{\sFmla{\False}{!A(t_1)}}
\RightLabel{$\TRule{\False}{\eq}$}
\UnaryInfC{\sFmla{\False}{!A(t_2)}}
\DisplayProof
\end{defish}
دھیان دیو کہ باقی سارے قاعدیاں دے برعکس، $\TRule{\True}{\eq}$ تے
$\TRule{\False}{\eq}$ لئی لازم اے کہ شاخ اُتے پہلے ای \emph{دو}
علامت لائے !!{formula}s موجود ہون، یعنی
$\sFmla{\True}{\eq[t_1][t_2]}$ تے $\sFmla{S}{!A(t_1)}$ دونیں۔

\begin{ex}
جے $s$ تے $t$ بند ٹرم نیں، تاں $\eq[s][t], !A(s)
\Proves !A(t)$:
\begin{oltableau}
  [\sFmla{\False}{\formula{A}(t)}, just = \TAss
    [\sFmla{\True}{\eq[s][t]}, just = \TAss
      [\sFmla{\True}{\formula{A}(s)}, just = \TAss
        [\sFmla{\True}{\formula{A}(t)}, just={\TRule{\True}{\eq}[2, 3]}, close]
      ]
    ]
  ]
\end{oltableau}
تساں ایس نوں اکو ہون والیاں چیزاں دی اک دوجے دی تھاں ورتوں دا اصول
(substitutability of identicals)، یا لائبنز دا قانون (Leibniz' Law)،
دے ناں نال جان دے ہوؤ گے۔

% PNBTAB-014: two explanatory identity-rule premise formulas are normalized; the three proof trees are unchanged.
!!^{tableau}s ایہ ثابت کردے نیں کہ $\eq$ متناظر اے، یعنی
$\eq[s_1][s_2]
\Proves \eq[s_2][s_1]$:
\begin{oltableau}
  [\sFmla{\False}{\eq[s_2][s_1]}, just = \TAss
    [\sFmla{\True}{\eq[s_1][s_2]}, just = \TAss
      [\sFmla{\True}{\eq[s_1][s_1]}, just = {$\eq$}
        [\sFmla{\True}{\eq[s_2][s_1]}, just = {\TRule{\True}{\eq}[2, 3]}, close]
      ]
    ]
  ]
\end{oltableau}
ایتھے سطر~$2$ پہلا لازمی !!{formula}~$\sFmla{\True}{\eq[s_1][s_2]}$
اے جیہڑا $\TRule{\True}{\eq}$ لئی چاہیدا اے۔ سطر~$3$ دوجا
لازمی فارمولا اے، جیہدا روپ $\sFmla{\True}{!A(s_1)}$ اے۔
$!A(x)$ نوں $\eq[x][s_1]$ سمجھو؛ فیر $!A(s_1)$،
$\eq[s_1][s_1]$ اے تے $!A(s_2)$، $\eq[s_2][s_1]$ اے۔

ایہ وی ثابت کردے نیں کہ $\eq$ متعدی اے، یعنی
$\eq[s_1][s_2],
\eq[s_2][s_3] \Proves \eq[s_1][s_3]$:
\begin{oltableau}
  [\sFmla{\False}{\eq[s_1][s_3]}, just = \TAss
    [\sFmla{\True}{\eq[s_1][s_2]}, just = \TAss
      [\sFmla{\True}{\eq[s_2][s_3]}, just = \TAss
        [\sFmla{\True}{\eq[s_1][s_3]}, just = {\TRule{\True}{\eq}[3, 2]}, close]
      ]
    ]
  ]
\end{oltableau}
ایس !!{tableau} وچ $\TRule{\True}{\eq}$ دا پہلا لازمی
!!{formula} سطر~$3$ دا $\sFmla{\True}{\eq[s_2][s_3]}$ اے
($s_2$،~$t_1$ دا کردار ادا کردا اے، تے $s_3$،~$t_2$ دا)۔
دوجا لازمی فارمولا، جیہدا روپ~$\sFmla{\True}{!A(s_2)}$ اے،
سطر~$2$ وچ اے۔ ایتھے $!A(x)$ نوں $\eq[s_1][x]$ سمجھو؛
ایس طرح $!A(s_2)$، $\eq[s_1][s_2]$ بندا اے (یعنی سطر~$2$)
تے $!A(s_3)$ نتیجے دا !!{formula}~$\eq[s_1][s_3]$ بندا اے۔
\end{ex}

\begin{prob}
ہِٹھ دتیاں لئی بند !!{tableau}s بناؤ:
\begin{enumerate}
\item $\sFmla{\False}{\lforall[x][\lforall[y][((x = y \land !A(x))
      \lif !A(y))]]}$
\item $\sFmla{\False}{\lexists[x][(!A(x) \land
    \lforall[y][(!A(y) \lif y = x)])]}$,\\
  $\sFmla{\True}{\lexists[x][!A(x)] \land
    \lforall[y][\lforall[z][((!A(y) \land !A(z)) \lif y = z)]]}$
\end{enumerate}
\end{prob}

\end{document}
